% Unit 064 English translation of chapter5.tex lines 9--140.
% Source: Methods of Algebra, Volume 2, commit
% 9a5803ff77dd3257484cb177f851a73770a59dd3 (CC BY 4.0).
% stable-unit-id: o014.aljabr2.chapter5.intro-filtration-grading
% Coverage: complete introduction to Chapter 5 and complete Section 5.1.

% segment-id: o014.li2.chapter5.intro-filtration-grading.h001
\chapter{Spectral Sequences}\label{sec:ss}
\hypertarget{o014.aljabr2.chapter5.intro-filtration-grading}{ }

% segment-id: o014.li2.chapter5.intro-filtration-grading.p001
Spectral sequences are classical algebraic tools that continually find new
uses, and their roots in topology run especially deep. Some textbooks and
lecture notes even use them as a basic tool of homological algebra to establish
some of the fundamental properties in \CHref{sec:Abel-cat} and
\CHref{sec:cplx}.

% segment-id: o014.li2.chapter5.intro-filtration-grading.p002
Spectral sequences were introduced by J.\ Leray in the 1940s and gradually
assumed their modern form through the work of J.-L.\ Koszul, H.\ Cartan, and
others. Their chief motivation was the study of the homology groups of fiber
bundles. In algebraic terms, the original problem amounts to understanding the
homology (or cohomology) of a chain complex (or complex) $X$. Let us take the
case of a complex as our example. Although $X$ itself may be difficult to
grasp, it can be studied through a suitable filtration
$\cdots\supset\mathrm{F}^pX\supset\mathrm{F}^{p+1}X\supset\cdots$.
The filtration $\mathrm{F}^\bullet X$ induces a filtration on every
$\Hm^n(X)$. Information about $\mathrm{F}^\bullet X$ and its subquotients can
then be used to approximate, step by step, each subquotient
$\gr^p\Hm^n(X)$ of the filtered cohomology. Spectral sequences are the art of
organizing such successive approximations.

% segment-id: o014.li2.chapter5.intro-filtration-grading.p003
A complex may also be generalized to a differential object $(X,d)$ in an
abelian category (Definition~\ref{def:differential-object}), leading to the
abstract definition of a spectral sequence in
Definition~\ref{def:SS-abstract}. A spectral sequence is a sequence of
differential objects $\mathscr{E}=(E_r,d_r)_r$ such that the cohomology of each
``page,'' $\Hm(E_r,d_r)$, is the next page $E_{r+1}$. Adding a grading and
allowing the differential $d$ to have a degree gives the notions of graded
spectral sequences $E_r^p$ and bigraded spectral sequences $E_r^{p,q}$
(Definition~\ref{def:graded-spectral-sequence}). Concretely, the differentials
in a bigraded spectral sequence have the form
\begin{gather*}
	d_r^{p, q}: E_r^{p, q} \to E_r^{p+r, q-r+1}, \\
	E_r = (E_r^{p, q})_{(p, q) \in \Z^2}, \quad d_r = (d_r^{p, q})_{(p, q) \in \Z^2}.
\end{gather*}
A filtered complex gives a filtered differential object and hence naturally
produces a bigraded spectral sequence. This is the form of spectral sequence
that occurs most often in applications. Every notion has a corresponding
version for chain complexes and their homology.

% segment-id: o014.li2.chapter5.intro-filtration-grading.p004
This chapter begins with W.\ Massey's theory of exact couples. Exact couples
are another device for producing spectral sequences, and their scope is wider
than that of filtered differential objects. After introducing filtrations and
graded structures in \S\ref{sec:grading-filtration}, we give the general
definition of a spectral sequence in \S\ref{sec:spectral-sequence}, together
with its graded and bigraded versions and such basic terminology as
degeneration, boundedness, and the limiting page. At the end of that section,
Remark~\ref{rem:edge-computing} explains how to read exact sequences from the
edges of a spectral sequence, an important technique.

% segment-id: o014.li2.chapter5.intro-filtration-grading.p005
Exact couples are defined in \S\ref{sec:exact-couples}. Next,
\S\ref{sec:filtered-diff-ss} explains how an exact couple produces a graded
spectral sequence from a filtered differential object, and
\S\ref{sec:filtered-cplx-ss} goes on to explain how a filtered complex produces
a bigraded spectral sequence. That section also treats the central notion of
convergence and the Classical Convergence Theorem~\ref{prop:classical-convergence}.
Some of the verifications are slightly technical, but there is no fundamental
difficulty. Readers may first work in a concrete abelian category such as
$\cate{Ab}$ to simplify the arguments.

% segment-id: o014.li2.chapter5.intro-filtration-grading.p006
A double complex gives rise to a filtered complex in two ways, according as
the total complex is filtered by the horizontal or the vertical coordinate.
The corresponding spectral sequences are denoted by
$\mathscr{E}_{\mathrm{I}}$ and $\mathscr{E}_{\mathrm{II}}$, respectively.
Their details are the subject of \S\ref{sec:double-cplx-ss}. We also give
standard applications, including the derivation of balanced bifunctors
(Example~\ref{eg:bifunctor-ss}), the spectral sequence of hyperderived functors
(Example~\ref{eg:hyperderived-ss}), and the Grothendieck spectral sequence
(Theorem~\ref{prop:Grothendieck-ss}). The Cartan--Eilenberg resolutions
introduced earlier (Theorem~\ref{prop:CE-resolution}) play a role here.

% segment-id: o014.li2.chapter5.intro-filtration-grading.p007
The Grothendieck spectral sequence concerns the derivation of a composite of
functors and is particularly common in geometry. As an example involving right
derived functors, consider left-exact additive functors between abelian
categories
$\mathcal{A}\xrightarrow{F}\mathcal{A}'\xrightarrow{F'}\mathcal{A}''$,
where $\mathcal{A}$ and $\mathcal{A}'$ have enough injective objects. If $F$
sends injective objects to $F'$-acyclic objects
(Convention~\ref{con:F-acyclic}), there is a canonical strongly convergent
bigraded spectral sequence
\[ E_2^{p, q} = (\mathrm{R}^p F') (\mathrm{R}^q F)(X) \Rightarrow \mathrm{R}^{p+q}(F' F)(X), \quad X \in \Obj(\mathcal{A}). \]
The derivation of composite functors can also be handled at the level of
derived categories; see
Theorem~\ref{prop:localization-triangulated-composite}. Besides being formally
elementary and easy to compute, the spectral-sequence version can provide
additional information in concrete situations, such as exact sequences in low
degrees.

% segment-id: o014.li2.chapter5.intro-filtration-grading.p008
The final section of the chapter, \S\ref{sec:multiplicative-SS}, briefly
introduces multiplicative structures on spectral sequences. This aspect has
received attention since the earliest history of spectral sequences.
Proposition~\ref{prop:dga-mult-ss} shows that a filtered differential graded
algebra gives rise to a spectral sequence with a canonical multiplicative
structure; a deeper treatment is left to specialized monographs.

% segment-id: o014.li2.chapter5.intro-filtration-grading.p009
As an exercise in applications, the Lyndon--Hochschild--Serre spectral
sequence, which is fundamental in group cohomology, can be obtained both as an
application of the Grothendieck spectral sequence and concretely from a
filtered complex. The latter approach also equips it with a canonical
multiplicative structure. This topic will be discussed in
\S\ref{sec:LHS-SS} in the future.

% segment-id: o014.li2.chapter5.intro-filtration-grading.note001
\begin{wenxintishi}
	This chapter requires only a knowledge of abelian categories and complexes.
	Apart from a few simple definitions, it does not depend on the material in
	\CHref{sec:triangulated-derived-cat}; the two chapters can be read
	independently. Since the spectral sequence of a filtered differential object
	can be written down directly, the theory of exact couples is not essential
	for that construction (see Proposition~\ref{prop:filtered-diff-E} and the
	discussion that follows). Readers may therefore consider temporarily
	skipping \S\ref{sec:exact-couples}. Moreover, among the later sections of
	this book, only \S\ref{sec:LHS-SS} requires a knowledge of spectral
	sequences.
\end{wenxintishi}

% segment-id: o014.li2.chapter5.intro-filtration-grading.h002
\section{Filtrations and Graded Structures}\label{sec:grading-filtration}

% segment-id: o014.li2.chapter5.intro-filtration-grading.p010
We encountered graded objects in Definition~\ref{def:graded-obj}. A
$\Z$-graded object in an arbitrary category $\mathcal{A}$ is simply called a
\emph{graded object}; by definition, these are the objects
$(X^p)_{p\in\Z}$ of $\mathcal{A}^{\Z}$. A $\Z^2$-graded object is simply
called a \emph{bigraded object}; by definition, these are the objects
$(X^{p,q})_{(p,q)\in\Z^2}$ of $\mathcal{A}^{\Z\times\Z}$. Morphisms between
graded (or bigraded) objects are written as $(f^p)_p$ (or
$(f^{p,q})_{p,q}$), as usual, and are composed degree by degree.
\index{graded object}

% segment-id: o014.li2.chapter5.intro-filtration-grading.p011
If $\mathcal{A}$ is an abelian category, then $\mathcal{A}^{\Z}$ and
$\mathcal{A}^{\Z\times\Z}$ are also abelian categories; kernels, cokernels,
and all other operations are computed degree by degree.

% segment-id: o014.li2.chapter5.intro-filtration-grading.p012
The category $\mathcal{A}^{\Z}$ has a shift autoequivalence $T$ determined by
$(TX)^p=X^{p+1}$ and $(Tf)^p=f^{p+1}$. More generally,
$\mathcal{A}^{\Z^n}$ has a family of shift autoequivalences
$T_1,\ldots,T_n$ corresponding to shifts along the $n$ coordinates. These
shifts commute strictly: $T_iT_j=T_jT_i$.

% segment-id: o014.li2.chapter5.intro-filtration-grading.p013
Within the framework of this chapter, graded structures arise primarily from
filtrations.

% segment-id: o014.li2.chapter5.intro-filtration-grading.d001
\begin{definition}\label{def:filtration}
	\index{filtration}
	\index{filtration!finite}
	\index[sym1]{FilpX@$\mathrm{F}^p X$, $\mathrm{F}_p X$}
	Let $X$ be an object of an additive category $\mathcal{A}$.
	\begin{itemize}
		\item A \emph{decreasing filtration} $\mathrm{F}^\bullet X$ is a
		sequence of subobjects
		\[ \cdots \supset \mathrm{F}^p X \supset \mathrm{F}^{p+1} X \supset \cdots \quad (p \in \Z). \]
		\item If there exist $a\leq b$ such that $\mathrm{F}^aX=X$ and
		$\mathrm{F}^bX=0$, the filtration is called \emph{finite}.
		\item Increasing filtrations $\mathrm{F}_\bullet X$ and their finiteness
		are defined similarly.
	\end{itemize}
\end{definition}

% segment-id: o014.li2.chapter5.intro-filtration-grading.p014
Decreasing and increasing filtrations can be distinguished by the position of
the index $p$. When no confusion is likely, both are simply called
filtrations. Decreasing filtrations are customary in the study of complexes
and cohomology, whereas increasing filtrations are customary for chain
complexes and homology. This chapter mainly uses decreasing filtrations.

% segment-id: o014.li2.chapter5.intro-filtration-grading.p015
All filtered objects $(X,\mathrm{F}^\bullet X)$ form a category
$\mathrm{Fil}^\bullet(\mathcal{A})$. A morphism
$f:(X,\mathrm{F}^\bullet X)\to(Y,\mathrm{F}^\bullet Y)$ is a morphism
$f:X\to Y$ that preserves the filtrations; that is, for every $p$, it
restricts to a morphism $\mathrm{F}^pX\to\mathrm{F}^pY$. The category
$\mathrm{Fil}_\bullet(\mathcal{A})$ is defined similarly.
\index[sym1]{FilA@$\mathrm{Fil}^\bullet(\mathcal{A})$, $\mathrm{Fil}_\bullet(\mathcal{A})$}

% segment-id: o014.li2.chapter5.intro-filtration-grading.p016
For every $X$, a decreasing filtration $\mathrm{F}^\bullet X$ (or increasing
filtration $\mathrm{F}_\bullet X$) gives a graded object
$(\mathrm{F}^pX)_{p\in\Z}$ (or $(\mathrm{F}_pX)_{p\in\Z}$). This defines a
functor $\mathrm{Fil}^\bullet(\mathcal{A})\to\mathcal{A}^{\Z}$ (or
$\mathrm{Fil}_\bullet(\mathcal{A})\to\mathcal{A}^{\Z}$), called the Rees
construction.

% segment-id: o014.li2.chapter5.intro-filtration-grading.p017
If $\mathcal{A}$ is an abelian category, there is another way to construct a
graded object from a filtration.

% segment-id: o014.li2.chapter5.intro-filtration-grading.d002
\begin{definition}
	\index[sym1]{grX@$\gr^p X$, $\gr_p X$}
	Let $\mathcal{A}$ be an abelian category, and consider a filtration
	$\mathrm{F}^\bullet X$ (or $\mathrm{F}_\bullet X$) on an object $X$.
	Define the graded object $\gr X$ by
	\[ \gr^p X := \mathrm{F}^p X / \mathrm{F}^{p+1} X \quad (\text{or} \quad \gr_p X := \mathrm{F}_p X / \mathrm{F}_{p-1} X). \]
	This defines a functor
	$\gr:\mathrm{Fil}^\bullet(\mathcal{A})\to\mathcal{A}^{\Z}$ (or
	$\gr:\mathrm{Fil}_\bullet(\mathcal{A})\to\mathcal{A}^{\Z}$).
\end{definition}

% segment-id: o014.li2.chapter5.intro-filtration-grading.p018
To extract information about $X$ from $(\gr^pX)_p$, we need at least
$\bigcup_p\mathrm{F}^pX=X$ and $\bigcap_p\mathrm{F}^pX=0$, using the notation
of Convention~\ref{con:increasing-union}. To make the first condition
meaningful and useful, we also want $\varinjlim_{p\to-\infty}$ to be exact in
$\mathcal{A}$, or else we may simply require that there be an $M$ such that
$\mathrm{F}^MX=X$. We begin with some related notions.

% segment-id: o014.li2.chapter5.intro-filtration-grading.d003
\begin{definition}\label{def:filtration-properties}
	\index{filtration!complete}
	\index{filtration!separated}
	\index{filtration!exhaustive}
	Let $\mathcal{A}$ be an abelian category and let $X$ be an object with a
	filtration $\mathrm{F}^\bullet X$.
	\begin{itemize}
		\item If $\bigcap_p\mathrm{F}^pX=0$, the filtration is called
		\emph{separated}.
		\item If $\bigcup_p\mathrm{F}^pX=X$ and either of the following
		conditions holds, the filtration is called \emph{exhaustive}:
		\begin{compactenum}[(a)]
		\item $\mathcal{A}$ has exact countable filtered direct limits
		$\varinjlim$; more
			precisely,
			$\varinjlim:\mathcal{A}^{(\Z_{\geq0},\leq)}\to\mathcal{A}$ exists
			and is exact\footnote{If $\mathcal{A}$ is a Grothendieck category,
			then (a) holds automatically.};
			\item there is an $M\in\Z$ such that $\mathrm{F}^MX=X$.
		\end{compactenum}
		\item If the family of canonical morphisms
		$X\twoheadrightarrow X/\mathrm{F}^pX$ induces an isomorphism
		$X\rightiso\varprojlim_pX/\mathrm{F}^pX$, the filtration is called
		\emph{complete}.
	\end{itemize}
	Increasing filtrations $\mathrm{F}_\bullet X$ are treated in the same way.
\end{definition}

% segment-id: o014.li2.chapter5.intro-filtration-grading.p019
The notion of completeness is inspired by the case of topological groups; see
\cite[\S 4.10]{Li1}.

% segment-id: o014.li2.chapter5.intro-filtration-grading.p020
A finite filtration is automatically separated, exhaustive, and complete;
this is the principal case needed in this chapter. Observe that if there is an
$N$ such that $\mathrm{F}^NX=0$, then $\mathrm{F}^\bullet X$ is separated and
complete. Moreover, since the kernel of
$X\to\varprojlim_pX/\mathrm{F}^pX$ is $\bigcap_p\mathrm{F}^pX$,
completeness implies separatedness.

% segment-id: o014.li2.chapter5.intro-filtration-grading.p021
The image of an exhaustive filtration is still exhaustive; this follows by
interpreting $\bigcup_p$ as $\sum_p$ and then applying
Lemma~\ref{prop:image-sum-preimage-intersection} (ii). Moreover, an exhaustive
filtration satisfies
\begin{equation}\label{eqn:exhaustion-intersection}\begin{gathered}
	\varinjlim_{p \to -\infty} \mathrm{F}^p X \rightiso \bigcup_p \mathrm{F}^p X = X, \\
	K = K \cap \left( \bigcup_p \mathrm{F}^p X \right) = \bigcup_p \left( K \cap \mathrm{F}^p X \right),
\end{gathered}\end{equation}
where $K$ is any subobject of $X$. Under condition (a) for an exhaustive
filtration, these assertions follow readily from
Proposition~\ref{prop:Grothendieck-cat-intersection}; under condition (b), they
are immediate.

% segment-id: o014.li2.chapter5.intro-filtration-grading.pr001
\begin{proposition}\label{prop:gr-isom}
	Let $\mathcal{A}$ be an abelian category, and let
	$f:(X,\mathrm{F}^\bullet X)\to(Y,\mathrm{F}^\bullet Y)$ be a morphism in
	$\mathrm{Fil}^\bullet(\mathcal{A})$.
	\begin{enumerate}[(i)]
		\item Suppose that $\mathrm{F}^\bullet X$ is separated and exhaustive. If
		$\gr(f)$ is a monomorphism, then $f$ is a monomorphism.
		\item Suppose that $\mathrm{F}^\bullet X$ is complete and exhaustive,
		while $\mathrm{F}^\bullet Y$ is separated and exhaustive. If $\gr(f)$ is
		an isomorphism, then $f$ is also an isomorphism, and in this case
		$\mathrm{F}^\bullet Y$ is complete as well.
	\end{enumerate}
	The corresponding statements hold for increasing filtrations.
\end{proposition}
% segment-id: o014.li2.chapter5.intro-filtration-grading.pf001
\begin{proof}
	For each $p\in\Z$, consider the following commutative diagram with exact
	rows:
% segment-id: o014.li2.chapter5.intro-filtration-grading.g001
	\[\begin{tikzcd}
		0 \arrow[r] & \mathrm{F}^{p+1} X \arrow[d, "f"'] \arrow[r] & \mathrm{F}^p X \arrow[r] \arrow[d, "f"] & \gr^p X \arrow[r] \arrow[d, "{\gr^p(f)}"] & 0 \\
		0 \arrow[r] & \mathrm{F}^{p+1} Y \arrow[r] & \mathrm{F}^p Y \arrow[r] & \gr^p Y \arrow[r] & 0
	\end{tikzcd}\]
	Write $K^p$ and $C^p$ for the kernel and cokernel of
	$\mathrm{F}^pX\xrightarrow{f}\mathrm{F}^pY$.

	For (i), apply Theorem~\ref{prop:snake-lemma} to the diagram above. This gives
	$K^{p+1}=K^p$. In particular,
	$K^p\subset\bigcap_{\ell\geq p}\mathrm{F}^\ell X=0$, so the restriction of
	$f$ to every $\mathrm{F}^pX$ is a monomorphism. To prove that $f$ itself is a
	monomorphism, it remains only to use
	\[ \Ker(f) \xlongequal{\because\; \text{\eqref{eqn:exhaustion-intersection}}} \bigcup_p \left( \Ker(f) \cap \mathrm{F}^p X \right) = 0. \]

	For (ii), the diagram gives both $K^{p+1}=K^p$ and
	$C^{p+1}\rightiso C^p$. For each $\ell>p$, consider the commutative diagram
	with exact rows
% segment-id: o014.li2.chapter5.intro-filtration-grading.g002
	\[\begin{tikzcd}
		0 \arrow[r] & \mathrm{F}^{\ell} X \arrow[r] \arrow[d, "f"] & \mathrm{F}^p X \arrow[r] \arrow[d, "f"] & \mathrm{F}^p X / \mathrm{F}^{\ell} X \arrow[r] \arrow[d] & 0 \\
		0 \arrow[r] & \mathrm{F}^{\ell} Y \arrow[r] & \mathrm{F}^p Y \arrow[r] & \mathrm{F}^p Y / \mathrm{F}^{\ell} Y \arrow[r] & 0
	\end{tikzcd}\]
	Another application of Theorem~\ref{prop:snake-lemma} gives the isomorphism
	induced by $f$
	\[ \mathrm{F}^p X / \mathrm{F}^\ell X \rightiso \mathrm{F}^p Y / \mathrm{F}^\ell Y. \]

	Taking $\varinjlim_{p:p<\ell}$ on both sides and using
	\eqref{eqn:exhaustion-intersection} gives
	$X/\mathrm{F}^\ell X\rightiso Y/\mathrm{F}^\ell Y$. Taking
	$\varprojlim_\ell$ now gives the commutative diagram
% segment-id: o014.li2.chapter5.intro-filtration-grading.g003
	\[\begin{tikzcd}
		X \arrow[r, "\sim"] \arrow[d, "f"'] & \varprojlim_\ell X/\mathrm{F}^\ell X \arrow[d, "\sim" sloped] \\
		Y \arrow[hookrightarrow, r] & \varprojlim_\ell Y/\mathrm{F}^\ell Y
	\end{tikzcd}\]
	It follows that both $f$ and
	$Y\to\varprojlim_\ell Y/\mathrm{F}^\ell Y$ are isomorphisms.
\end{proof}
